\section{Intermediate powers of bias}
\label{sec:intermediate-powers-of-bias}

The preceding reductions give discrepancy at small and unequal power
widths. We now compare available bias across those budgets and budgets that
are linear in \(n\). Two limiting exponents record the transition. The goal
is to exclude intermediate values, leaving the possible jump at a linear
budget for Section~\ref{sec:jump-to-large-bias-only-at-linear-budget}, after
the cluster exclusion.

\begin{definition}[Limiting bias exponents]\label{def:limiting-bias-exponents}
Stabilize the absolute-bias properties at all rational budgets and exponents used below. In a specified orientation \(X,Y\), let \(H(x,y)\) be the infimum of rational \(h>0\) for which bias of magnitude at least \(n^{-h}\) is available at widths \(n^x,n^y\), where \(0<x,y<1\). Set the infimum to \(+\infty\) if there is no such exponent. If \(h<H(x,y)\), stabilization makes this witness property eventually absent: every pair of laws within these budgets then has bias less than \(n^{-h}\). Let \(H_L(x,\alpha)\) denote the corresponding infimum with second budget \(\alpha n\), for \(\alpha>0\).

Both functions are nonincreasing in their width parameters, since enlarging a budget preserves availability. Taking limits through rational arguments, put
\[
 H^\dagger=\lim_{x\downarrow0}\lim_{y\uparrow1}H(x,y),\qquad
 H_L^\dagger=\lim_{x,\alpha\downarrow0}H_L(x,\alpha).
\]

\end{definition}

\begin{proposition}[Exclusion of intermediate bias]\label{prop:intermediate-bias-exclusion}
In the stabilized contradiction sequence, neither of the following situations can occur:
\begin{itemize}
\item some \(H^\dagger<1\);
\item both \(H^\dagger\ge1\), and some \(0<H_L^\dagger<1\).
\end{itemize}
\end{proposition}

\begin{proof}\leavevmode\par
\paragraph{Parameter selection.}
Here are parameter choices and what is then available, using the orientation in question. We arrange \(0<x_s<x_d<.1,\ u=x_s/2\), exponents \(0<h_-<h_+<1\), \(x_d<(1-h_+)/10\), and small \(\sigma,\chi>0\) with
\[
 \sigma<x_s/10,\qquad
 \chi<\min(x_s,h_-,1-h_+)/100,\qquad h_+-h_-<\chi/10.
\]
Write \(a_*=\tfrac12 n^{-h_+},\ b_*=n^{-h_-}\). At shallow budgets \(n^{x_s},S_s\) absolute bias \(\ge2a_*\) is available; at deep budgets \(n^{x_d},S_d\) all densities are \(1/2\pm b_*\) (meaning within the indicated error). Further choices:
\paragraph{Sublinear case.} \(S_s=n^{y_s}, S_d=n^{y_d}\), with \(y_s<1-\sigma<y_d\), \(\chi<\sigma/10\); use \(m=\lfloor n^{y_m}\rfloor\) special bits with \(y_s<y_m<1-\sigma\);
\paragraph{Linear case.} \(S_s=\alpha_s n,\ S_d=\alpha_d n\), with \(100\alpha_s<\alpha_d<.01,\ \sigma<\chi/10\); use \(m=\lfloor\alpha_d n/10\rfloor\). Also densities against any law of width \(O(1)\) on \(Y\), for all first laws up to width \(n^{x_d}\), are \(1/2\pm o(a_*)\).

The two cases use different orders of parameter selection: the inequalities specific to either case are imposed only in that case.

For the sublinear case, the thresholds we need are
\[
 h_-<H(x_d,y_d)\le H(x_s,y_s)<h_+.
\]
Choose a small rational interval \([x_1,x_2]\) within the first-budget range of Corollary~\ref{cor:asymmetric-discrepancy}, also small relative to the margin between \(H^\dagger\) and 1. Choose a rational \(y\)-interval near, but bounded away from, 1. We may arrange that \(H\) is bounded above by a fixed number below 1 throughout the resulting rectangle, that \(x_2\) is less than a small constant times this margin, and that \(1-y\ll x_1\). The value at the deep corner is positive by the asymmetric discrepancy corollary; monotonicity gives a positive lower bound throughout the rectangle.

Fix \(\chi\) using these positive margins, including \(\chi<(1-y)/10\) throughout the rectangle. Subdivide its ascending diagonal into sufficiently many equal intervals. The total decrease of \(H\) is bounded, so one adjacent pair has decrease less than \(\chi/20\). Use its endpoints as the shallow and deep budgets, and place \(h_-,h_+\) strictly below and above their respective \(H\)-values, with \(h_+-h_-<\chi/10\). We can then choose \(1-\sigma\) strictly between \(y_s,y_d\), and choose \(y_m\) between \(y_s\) and \(1-\sigma\). These strict choices provide the required width slack.

In the linear case, first choose fixed positive margins around \(H_L^\dagger\in(0,1)\). The nonincreasing function
\(x\mapsto\lim_{\alpha\downarrow0}H_L(x,\alpha)\)
tends to \(H_L^\dagger\) as \(x\downarrow0\). Choose a small \(x\)-interval on which these values stay within the fixed margins. Since \(H^\dagger\ge1\), the interval can also be chosen so that a sublinear discrepancy bound with exponent \(h_{\rm broad}>h_+\) holds there. In particular, tests against laws of width \(O(1)\) have error \(n^{-h_{\rm broad}}=o(a_*)\).

Fix \(\chi\) using this interval and the margins. As above, monotonicity gives \(x_s<x_d\) for which the decrease of the limiting linear exponent is less than \(\chi/30\). Choose \(\alpha_d\) sufficiently small, and then \(\alpha_s<\alpha_d/100\) sufficiently small, to obtain
\[
 h_-<H_L(x_d,\alpha_d)\le H_L(x_s,\alpha_s)<h_+.
\]
Here \(\sigma\) may be chosen arbitrarily small after \(\chi\), so the linear-case conditions on it are compatible.

In either case, availability at the shallow threshold permits one color to be fixed along a subsequence. Restrict each first law to labels of degree at least \(1/2+a_*\) into its second law. Since the original surplus is at least \(2a_*\), the retained mass is at least a constant times \(a_*\); this costs \(O(\log n)\) width. Lemma~\ref{lem:balanced-profiles} now gives a balanced finite mixture of patches \((\mu_i,\nu_i)\). Each first support has the stated minimum degree, the shallow widths have only this logarithmic extra cost, and the deep discrepancy bounds hold on all retained labels.

\subsection{Core and read bound for center IDs}
\label{sec:core-and-read-bound-for-center-ids}

Write a cube word as a special word \(z\in Q_m\) and a residual word in \(Q_{n-m}\). A center ID records a slice \(z\), a residual location and a level. We first construct and then fix a coloring-independent map from every cube site, of either parity, to an ID in its own slice at residual distance at most \(r=\lfloor n^\sigma\rfloor\). Randomness is used here only to prove existence of this map. The anchor variables introduced later are fresh, independent draws.

The map will have the following properties:
\begin{itemize}
\item At an odd row, the adjacent even sites use at most \(T+m\) distinct IDs, where \(T=O(n^{1-\sigma+\varepsilon})\) and \(\varepsilon>0\) can be arbitrarily small.
\item At an even row \(v=(z,v')\), the IDs observed in all neighboring odd stars have a common core of size at most \(n^\chi\), including the ID of \(v\). Every other ID occurs in at most \(r+3\) of these stars.
\end{itemize}

\paragraph{The adapted height construction.}
Let \(V_j\) be the volume of a radius-\(j\) ball in the residual cube, and put \(V=V_r\). Independently at every slice/location/level, put a prospective center with probability \(n^{10}/V\), and activate it with probability \(n^{b_0-10}\). Here \(b_0>0\) will be small. All prospective centers in the same-slice radius-\(r\) ball are eligible. Besides the event that this ball contains no active center, use the following crowd tests:
\[
\begin{array}{c|c|c}
 \text{active centers counted}&\text{mean}&\text{upper threshold}\\ \hline
 \text{same slice, radius }r&n^{b_0}&n^{\chi/2}\\
 \text{adjacent slices, radius }r-1&O(mn^{b_0}r/n)&n^{\chi/2}\\
 \text{same slice, radius }r+1&O(n^{b_0}n/r)&n^{1-\sigma+\varepsilon}
\end{array}
\]
Require these bounds at the levels within 2 of the level being tested, clipping that window at its endpoints. The mean estimates use \(V_{r-1}/V=O(r/n)\) and \(V_{r+1}/V=O(n/r)\). Choose \(0<\varepsilon<\sigma\) within the filter-budget slack below, and then \(b_0<\varepsilon\) sufficiently small in terms of \(\chi\). In the linear case also require \(b_0+\sigma<\chi/2\); in the sublinear case \(y_m+\sigma<1\) provides the corresponding slack. Thus each crowd threshold exceeds its mean by a fixed power of \(n\).

We adapt Steps 2--4 of the proof of Lemma~\ref{lem:height-selection}, using paths with step bound 2 in the full cube. The position counts in all eligibility balls are at least \(n^{10}/2\) with high probability. The hole and crowd probabilities, including the bounded level windows, have the same base estimates as in that proof. The following overlap check permits the paths to cross slices.

A child domain of metric radius \(R'\) consults slices within \(O(R')\) special-coordinate distance and residual locations within \(r+O(R')\). For starts separated by \(KR'\), whose level ranges overlap, either the consulted slice ranges are disjoint or the residual separation is \(\Omega(KR')\). If the residual balls also intersect, their centers are at distance at most \(2r+O(R')\), so \(R'=O(r)\) for large fixed \(K\). The shell and hypergeometric estimates from the height proof then bound their residual overlap fraction by
\(\exp(-\Omega(KR'\log n))\).
Enlarging the residual balls and enumerating the consulted slices costs only \(\exp(O(R'\log n))\). Taking \(K\) large therefore leaves pair-level overlap volume at most \(Ve^{-\Omega(R')}\).

Use degraded thresholds at successive scales for every crowd count in the table and for the eligible-size count. Delete overlaps to obtain private child regions. Deletion preserves holes; its loss from each size or crowd count fits the corresponding threshold gap unless an overlap contains too many prospective or active centers. With \(q\) separated child domains, the binomial bound for each such exception is at most
\[
 \exp(-\Omega(n^{\chi/2}R'/q)).
\]
The minimum crowd exponent is \(\chi/2\). Denote the height proof's scale-spacing exponent by \(\sigma_h\), to distinguish it from the present radius exponent. Choose its \(a,\zeta,1-\theta,\sigma_h\) small enough that
\[
 0<\sigma_h<\zeta,\qquad
 b_0>a>\zeta+\sigma_h+(1-\theta),\qquad
 \chi/2>a+4\sigma_h.
\]
The base bound, the overlap bound and independence on private regions now give the same scale induction as in Lemma~\ref{lem:height-selection}. This argument uses the three displayed crowd tests; it does not require the original lemma's full-cube radius-\((r+2)\) crowd test. It gives good heights differing by at most one at full-cube distance at most 2. Choose a covered active center at every such height and fix the resulting map.

\paragraph{The core and its read multiplicity.}
At an odd star, the residual-flip neighbors use at most \(O(n^{1-\sigma+\varepsilon})\) IDs by the radius-\((r+1)\) crowd bound and the bounded level window. Its special-flip neighbors use at most \(m\) more. At \(v\), put in the core every observed ID in the same-slice radius-\(r\) ball or in an adjacent-slice radius-\((r-1)\) ball, at the relevant heights within 1. The two smaller crowd bounds make the core size at most \(n^\chi\).

To verify the read bound, let \(\Delta(c,v')\) be the set of residual coordinates where a center \(c\) differs from \(v'\). Suppose a same-slice ID outside the core is seen at the residual-flip neighbor \(b_i\) through an even neighbor \(v'^{\,i,j}\). The case \(j=i\) returns to \(v'\) and would put the ID in the core. Otherwise
\[
 d_H(c,v'^{\,i,j})=d_H(c,v')+2
   -2\bigl(1_{i\in\Delta(c,v')}+1_{j\in\Delta(c,v')}\bigr).
\]
The left side is at most \(r\), while \(d_H(c,v')>r\). Hence \(i\in\Delta(c,v')\), and \(d_H(c,v')\le r+2\). Such an ID is therefore seen at no more than \(r+2\) residual-flip neighbors. For an adjacent-slice ID outside the core, the one-flip identity gives \(i\in\Delta(c,v')\) and \(d_H(c,v')\le r+1\). At special-flip neighbors, same-slice IDs are already in the core, an adjacent slice contributes only at its own special flip, and a distance-two slice contributes at most twice. These bounds give \(r+3\).

Finally choose \(\varepsilon\) so that
\[
 S_s+O(n^u)+\log(1/.49)(T+m)<S_d.
\]
In the sublinear case both \(T\) and \(m\) have exponent below \(y_d\); in the linear case \(T=o(n)\) and \(m=\lfloor\alpha_dn/10\rfloor\), while \(100\alpha_s<\alpha_d\). Thus this choice is possible.

\paragraph{Tags and masks.}
Assign independent tags \(i_z\) from the balanced mixture, one per special word. Their averaged first and second laws have pointwise normalized loads \(O(1)\) with high probability. In the \(n\)-th moment, a repeated word costs a fraction at most \(n2^{-m}\), multiplied by its normalized atom cap. This product tends to zero: \(m\gg S_s\) in the sublinear case, and \(m\log2> S_s\) with linear slack in the linear case. The first-law cap is smaller still.

In the linear case we also choose the tags so that, for every \(z\), except on \(\mu_{i_z}\)-mass \(\exp(-\Omega(n^u))\),
\[
 \sum_{j\le m}\bigl(d_G(w,\nu_{i_{z^j}})-1/2\bigr)
       \ge -.05a_*n .
\]
Here \(z^j\) is obtained by flipping bit \(j\). To justify this choice, first use the raw product law of the tags and draw \(w\sim\mu_{i_z}\). The broad-test bound gives degree \(1/2\pm o(a_*)\) into the tag-average second law outside \(\exp(-\Omega(n^u))\) first-law mass. Indeed, conditioning on a larger exceptional set would stay inside the first deep budget. Deep discrepancy also gives individual degrees \(1/2\pm2b_*\) except with the same exponential scale of joint \((w,i_{z^j})\)-probability. By Markov, outside an exponentially small set of \(w\)'s, the conditional probability of such an irregular neighboring tag is exponentially small.

Clip every summand to \([-2b_*,2b_*]\). Conditional on \(w\), these clipped summands depend on independent adjacent tags. Their downward deviation of order \(na_*\) has probability
\(\exp(-\Omega(na_*^2/b_*^2))\).
Adding the clipping exceptions proves a joint tail \(\exp(-\Omega(n^u))\). If \(E_z(w)\) denotes failure of the displayed sum, choose a small fixed \(c>0\) and put
\[
 B_z=\{\mu_{i_z}(E_z)>e^{-c n^u}\}.
\]
Markov gives \(\Pr(B_z)=\exp(-\Omega(n^u))\). The event \(B_z\) reads only the tag at \(z\) and its neighbors, so the dependency degree is polynomial. Lemma~\ref{lem:conditional-avoidance} allows all \(B_z\)'s to be avoided with exponentially small charges. The preceding tag-load moment still applies: removing constraints touching at most \(n\) separated tag neighborhoods has charge sum \(\operatorname{poly}(n)e^{-c' n^u}=o(1)\). Fix tags satisfying these conditions and the load bounds.

Each ID \(c\) in slice \(z\) now carries an independent anchor \(W_c\sim\mu_{i_z}\). At an odd row \(b\), let \(\nu=\nu_{i_{z(b)}}\). Choose a mask of \(\nu\)-mass at least \(\tfrac12e^{-n^u}\), normalize its restriction, and require hits of the anchors at every distinct neighboring ID. On a nonempty filter let \(p_b\) be the normalized result; otherwise use the initial masked law as an auxiliary fallback.

Choose the mask distribution separately at each row by convex separation, before drawing the anchors. For any nonnegative price, the labels costing at most \(1+e^{-n^u}\) times its \(\nu\)-mean have mass at least \(\tfrac12e^{-n^u}\), and hence form an allowed mask. Every output from that mask has the same price bound, including the fallback. Separation therefore supplies a distribution of masks for which
\[
 R_b:=\mathbb E p_b\le(1+e^{-n^u})\nu.
\]
The expectation here averages all anchors used at \(b\) and its mask, before any further conditioning. Masks are drawn independently across rows. Since \(R_b\) and \(\nu\) are probability laws, their total-variation distance is exponentially small. This is an unconditional mean identity; it is not asserted after a core has been fixed.

\subsection{Gains after conditioning the core}
\label{sec:gains-after-conditioning-the-core}

All probabilities in this subsection use independent anchors and masks, with the tags and ID map fixed. For an even role \(v\), write \(W_*=W_{c(v)}\). At each \(b\sim v\), let \(q_b\) be the fraction retained by imposing the hit of \(W_*\) after all other hits in that filter. We will prove, with failure probability \(\exp(-\Omega(n^u))\) per even star, that all these filters have width less than \(S_d\) and
\begin{equation}
 \sum_{b\sim v}\log q_b\ge-n\log2+c_4na_* ,
 \label{eq:source-3}
\end{equation}
where \(c_4>0\) is a small absolute constant.

\paragraph{The required regularity tests.}
Fix an ordering of all ID names. For each \(b\sim v\), order its IDs outside the core first, then its core IDs other than \(c(v)\), and finally \(c(v)\), keeping the fixed order inside the first two blocks. We use this order with the masked starting law and with the unmasked starting law. We also use the order of the core IDs other than \(c(v)\), starting from the unmasked law. Include all prefixes of these three orders. The masked prefixes give the outer filter, the deletion filter and the target-last fraction; the unmasked core prefixes are the laws used in the covariance telescoping argument below. There are \(O(n(T+m))\) such prefix tests per star. The degree and signed-density tests below are applied at these same prefixes, so their number is polynomial as well. No assertion over all permutations or all subsets of IDs is needed.

At a fixed prefix whose preceding hits each retain at least \(.49\), the current second-law width is at most
\[
 S_s+O(n^u)+\log(1/.49)(T+m)<S_d,
\]
with room for the further \(O(n^u)\) restrictions below. If the next hit fraction differs from \(1/2\) by more than \(2b_*\) on anchor mass at least \(e^{-O(n^u)}\), condition its first law on one signed exceptional set. Its width is at most
\[
 n^{x_s}+O(n^u)+O(\log n)<n^{x_d},
\]
and its density against the current second law contradicts deep discrepancy. Thus each next hit fraction is \(1/2\pm2b_*\), except with probability \(\exp(-\Omega(n^u))\). Conditional application at each prefix and a union over the stated family prove all required sequential tests. We enlarge the constant implicit in \(T\) if necessary.

We will also use the reverse form of this reasoning: for any first law within the deep budget, an exceptional set of second labels of mass at least \(e^{-n^u}\) on which its degree differs by more than \(2b_*\) could be normalized and used as a second test law. The available width room excludes this. In particular, each actual \(q_b=1/2\pm2b_*\) except with the stated tail. The remaining issue is a positive sum of their surpluses.

For the fixed even role \(v\), let \(C_v\) be the core specified above, common to all its neighboring stars. If \(I_b\) is the ID set observed at \(b\sim v\), put
\[
 D_b=I_b\cap C_v,\qquad O_b=I_b\setminus C_v.
\]
Condition on the entire anchor vector \(W_{C_v}=(W_c)_{c\in C_v}\), keeping the tags and the deterministic ID map fixed. At a given \(b\), write \(D=D_b\), \(O=O_b\), and \(k=|D|\le n^\chi\); the target ID \(c(v)\) belongs to \(D\). Let \(J_-,J_D\) be the sets of second labels hitting \(D\) except the target, or all of \(D\), respectively. Let \(P_O\) be the masked law filtered just by outer hits, with a width-preserving fallback (initial masked law) if their retained mass is below \(.49^{|O|}\). Write \(Q=\mathbb E P_O\), averaging the outer anchors and this row's independent mask; thus \(Q\) is independent of \(W_{C_v}\). Core coordinates outside \(I_b\) do not enter this filter. Clip \(q_b\) to \([1/2-2b_*,1/2+2b_*]\), with a fixed local handling of invalidity there (using only this filter's inputs), calling it \(\widehat q_b\).

Conditional Markov applied to the outer-then-core tests shows that, outside a set of core histories of mass \(\exp(-\Omega(n^u))\), their conditional failure probability is also \(\exp(-\Omega(n^u))\). On these histories, outside that conditional exception,
\[
 P_O(J_-)=2^{1-k}(1+O(kb_*)),\qquad
 P_O(J_D)=P_O(J_-)\widehat q_b.
\]
Consequently
\[
 d_G(W_*;Q|_{J_-})
 =\frac{\mathbb E[P_O(J_-)\widehat q_b\mid W_{C_v}]}
        {\mathbb E[P_O(J_-)\mid W_{C_v}]}
   +O\bigl(2^k e^{-\Omega(n^u)}\bigr).
\]
The fraction is a weighted mean of \(\widehat q_b\). Its weight differs from the constant \(2^{1-k}\) by relative \(O(kb_*)\), while \(\widehat q_b\) has oscillation \(O(b_*)\). Removing this weight therefore costs \(O(kb_*^2)\), rather than \(O(kb_*)\). Since \(k\le n^\chi=o(n^u)\), the exceptional error remains exponentially small. We obtain
\[
 \mathbb E[\widehat q_b\mid W_{C_v}]
 =d_G(W_*;Q|_{J_-})+O(kb_*^2)+\exp(-\Omega(n^u))O(1).
\]
Here \(d_G(w;\lambda)\) is the degree of one label into a probability law, and every restriction bar denotes normalization.

\paragraph{Erasing the core from the mean filter.}
The law \(Q\) averages only the outer anchors and the mask. To compare it with \(\nu\), temporarily average the core as well, returning to the unconditional law \(R_b\) above. Independence of distinct core anchors gives the exact identity
\[
 2^k\mathbb E[P_O(y)1_{J_D}(y)]=a_D(y)Q(y),\qquad
 a_D(y)=2^k\prod_{c\in D}d_G(\mu_{i_{z(c)}},y).
\]
On regular histories the full filter is \(P_O|_{J_D}\), and its normalizing multiplier is \(2^k(1+O(kb_*))\). Averaging this bounded multiplier pointwise yields a function \(\theta\) with \(|\theta|\le Ckb_*\) such that
\[
 R_b=(1+\theta)a_DQ+e,\qquad
 \|e\|_1=\exp(-\Omega(n^u))O(1).
\]
To see the error estimate, extend the good-history multiplier divided by \(2^k\) by 1 on bad histories. Replacing the actual kernel there costs at most their probability times \(1+2^k\), which preserves the exponential scale.

The fallback in \(P_O\) gives it, and therefore \(Q\), width at most
\(S_s+O(n^u)+\log(1/.49)|O|\).
The reverse degree tests just proved apply to \(Q\). Outside a \(Q\)-set of mass \(\exp(-\Omega(n^u))\), every doubled degree factor \(2d_G(\mu_{i_{z(c)}},y)\) is \(1+O(b_*)\), so \(a_D=1+O(kb_*)\). On the exceptional set use \(a_D\le2^k\). Since \(kb_*=o(1)\), we can invert the resulting multiplier and use the exponentially small distance between \(R_b\) and \(\nu\). This gives
\[
 Q=(1+\psi)\nu+\mathrm{err},\qquad
 |\psi|\le O(kb_*),\qquad
 \|\mathrm{err}\|_1=\exp(-\Omega(n^u))O(1).
\]
All measures and multipliers in this identity are independent of the core data.

Now fix the core except for \(W_*\). The unmasked regularity tests give \(\nu(J_-)\ge.49^k\) outside an exponentially small exception. Restricting the last identity to \(J_-\) magnifies its total-mass error by at most \(O(.49^{-k})\), still negligible. Put \(\lambda=\nu|_{J_-}\). The normalized multiplier changes \(\lambda\) by a signed density
\[
 s=\frac{\psi-\mathbb E_\lambda\psi}{1+\mathbb E_\lambda\psi},\qquad
 \mathbb E_\lambda s=0,\qquad |s|\le Ckb_*.
\]
Thus \(\tau=(1+s/(Ckb_*))\lambda\) is a probability law whose width exceeds that of \(\lambda\) by at most \(\log2\). The target anchor \(W_*\) is independent of these two laws. The individual-degree test gives \(d_G(W_*;\tau)-d_G(W_*;\lambda)=O(b_*)\), except with the stated exponential tail. Rescaling back to the actual modulation gives
\[
 d_G(W_*;(1+s)\lambda)-d_G(W_*;\lambda)
 =Ckb_*\bigl(d_G(W_*;\tau)-d_G(W_*;\lambda)\bigr)
 =O(kb_*^2).
\]
Therefore replacing \(Q|_{J_-}\) by \(\nu|_{J_-}\) in the conditional mean costs \(O(kb_*^2)\).

\paragraph{The effect of each core hit.}
It remains to compare \(d_G(W_*;\nu|_{J_-})\) with \(d_G(W_*;\nu)\). At a prefix of the unmasked core order, let \(\lambda\) be the current law, let \(w\) be the next anchor, and write \(x=W_*\). These two anchors are independent of \(\lambda\) and each other. With \(g_w(y)=1[w\sim_Gy]\), the degree change on imposing the next hit is
\[
 d_G(x;\lambda|_{N_G(w)})-d_G(x;\lambda)
 =\frac{\operatorname{cov}_\lambda(g_w,g_x)}{d_G(w;\lambda)}.
\]
The denominator is at least \(.49\) on regularity. We claim that the covariance has magnitude at most \(a_*n^{-2\chi}\), except with probability \(\exp(-\Omega(n^u))\).

Suppose instead that its bad-pair probability is at least \(e^{-n^u}\). Choose one sign and an \(x\)-set \(X_0\) of mass at least \(\rho=e^{-O(n^u)}\) such that, for every \(x\in X_0\), the \(w\)-set \(B_x\) with that signed violation has mass at least \(\rho\). Draw \(t=\lceil n^{8\chi}\rceil\) independent \(w_j\)'s from the anchor law conditioned on \(B_x\). This conditioned law still fits the first deep budget. Its mean adjacency indicator is \(1/2\pm2b_*\) outside a \(\lambda\)-set of exponentially small mass, by the reverse degree test. Hence, for
\[
 S=\sum_{j\le t}(g_{w_j}-\mathbb E_\lambda g_{w_j}),
\]
independence gives
\[
 \mathbb E\|S\|_{L^2(\lambda)}^2
       =O(t+t^2b_*^2)=O(t),
\]
since \(tb_*^2=o(1)\). For an absolute \(K\), the tuple event \(\mathcal G=\{\|S\|_{L^1(\lambda)}\le K\sqrt t\}\) thus has conditional probability at least a fixed \(c>0\).

The event \(\mathcal G\) depends on the tuple and the fixed law \(\lambda\), not on \(x\). Returning to the raw independent tuple law gives
\[
 \int_{X_0}\Pr_{\rm raw}
     (\mathcal G,\ w_j\in B_x\text{ for all }j)\,d\mu(x)
       \ge c\rho^{t+1}.
\]
Here \(\mu\) denotes the target anchor law. Fubini therefore supplies one fixed norm-good tuple whose simultaneous signed-violation set of \(x\)'s has mass at least \(c\rho^{t+1}=e^{-O(tn^u)}\). Restricting \(\mu\) to that set costs \(O(n^{u+8\chi})\) width, which fits the first deep budget because \(u+8\chi<x_s\).

For this fixed tuple, \(\mathbb E_\lambda S=0\), so the laws
\[
 \lambda_\pm=\frac{\sqrt t+S_\pm}{Z}\lambda,
 \qquad Z=\sqrt t+\mathbb E_\lambda S_+
           =\sqrt t+\mathbb E_\lambda S_-=O(\sqrt t)
\]
have equal normalizers. Here \(S_\pm=\max(\pm S,0)\). Their multipliers are \(O(\sqrt t)\), costing only \(O(\log n)\) second width. Deep discrepancy, with the target law restricted to the simultaneous violation set, bounds the average over that law of \(\mathbb E_\lambda(g_xS)\) in magnitude by \(O(b_*\sqrt t)\). The simultaneous violations force it to exceed \(ta_*n^{-2\chi}\) in the chosen sign. This is impossible, since
\[
 \frac{\sqrt t\,a_*n^{-2\chi}}{b_*}
       \asymp n^{2\chi-(h_+-h_-)}\longrightarrow\infty.
\]
This proves the covariance tail.

There are at most \(k\le n^\chi\) core hits. Their total degree shift is \(O(ka_*n^{-2\chi})=o(a_*)\). Also \(kb_*^2=o(a_*)\). Combining this with the erasure comparison, and taking the union over \(b\sim v\), gives outside exponentially rare core data
\[
 \mathbb E[\widehat q_b\mid W_{C_v}]
       =d_G(W_*;\nu_{i_{z(b)}})+o(a_*).
\]

\paragraph{Concentration with the core fixed.}
The comparisons above hold for the fixed tags. Every ordinary neighbor uses the same slice as \(v\), so its conditional mean surplus is at least \(a_*+o(a_*)\). Their sum is at least \((1-o(1))(n-m)a_*\). In the linear case the additional tag condition bounds the sum over special neighbors below by \(-.06na_*\). In the sublinear case the bound \(mb_*=o(na_*)\) gives the required estimate without that condition.

After fixing \(W_{C_v}\), the remaining inputs are independent outside-core anchors and independent row masks. Each clipped fraction reads one row mask, and each remaining anchor is read at most \(r+3\) times. Exponential Markov and the product-space H\"older inequality~\cite{finner1992} therefore bound a downward deviation of \(.1na_*\) by
\[
 \exp\!\left(-\Omega\!\left(\frac{na_*^2}{(r+3)b_*^2}\right)\right).
\]
Here the product-space inequality bounds the expectation of a product of nonnegative functions by the product of their \(d\)-norms when each independent input is read at most \(d\) times. We apply it with \(d=r+3\) in this raw conditional experiment, before avoidance is imposed. The parameter choices give
\[
 \frac{na_*^2}{(r+3)b_*^2}
      =n^{1-\sigma-2(h_+-h_-)+o(1)}\gg n^u.
\]
Thus the tail has the required strength. On individual regularity, the clipped fractions equal the actual \(q_b\)'s. Finally
\[
 \log(\tfrac12+t)=-\log2+2t+O(t^2),\qquad |t|\le2b_*.
\]
The quadratic loss is \(O(nb_*^2)=o(na_*)\), so the positive surplus proves \eqref{eq:source-3}.

\subsection{Taking the assignment}
\label{sec:taking-the-assignment-intermediate}

\paragraph{Anchor avoidance and the predictive tests.}
At every even star require the neighboring filter tests just enumerated and the log-gain bound \eqref{eq:source-3}. For a target \(v\), fix all anchors except \(W_*\) and all masks. For candidate \(W_*=x\), define the data sublikelihood
\[
 F_x(\mathbf y)=1_{\mathcal V}\prod_{b\sim v}p_b(y_b),
\]
where \(\mathcal V\) is this star's filter-validity and log-gain event, recomputed with \(x\). Let \(P^-\) be the product of the deletion kernels, each obtained by normalizing the same filter without the target ID. A zero deletion denominator uses a fixed fallback independent of \(x\). On \(\mathcal V\), \eqref{eq:source-3} gives
\[
 F_x\le2^n e^{-c_4na_*}P^-.
\]
Put \(M_v=\int F_x\,d\mu_{i_z}(x)\), where \(z=z(v)\). Define predictive failure by
\(M_v=0\) or \(M_v<e^{-c_4na_*/4}P^-\).
Lemma~\ref{lem:gated-posterior} bounds its true-gated probability in the raw anchor-and-product-data experiment by \(e^{-c_4na_*/4}\). Markov then shows that the anchor event
\[
 1_{\mathcal V}\Pr_{\prod_{b\sim v}p_b}
          (\text{predictive failure})>e^{-c_4na_*/8}
\]
has probability at most \(e^{-c_4na_*/8}\).

Each star event reads IDs within residual distance \(r+O(1)\) and bounded special-coordinate distance. Its dependency degree, and the number of events touching any one local list, are at most
\(\operatorname{poly}(n)e^{O(r\log n)}\).
Since \(\sigma<u\), charges \(e^{-c n^u}\), with small fixed \(c>0\), meet Lemma~\ref{lem:conditional-avoidance} for the star and predictive events. Moreover,
\[
 \operatorname{poly}(n)e^{O(r\log n)}e^{-c n^u}=o(1),
\]
so removing the touching constraints for any polynomial number of queried local lists costs \(1+o(1)\).

\paragraph{Odd column loads and injection.}
Fix a second label \(y\) and use normalized weights \(Np_b(y)\) in Lemma~\ref{lem:scattered-moments}. A row has cap \(e^{S_d}\) on success. Declare rows near when their special words are at bounded distance and their residual words at distance \(O(r)\), with constants large enough to contain the declared input lists. The near fraction is at most \(2^{-n}\exp(O(r\log n))\), and
\[
 \exp(O(r\log n)+S_d-n\log2)=o(1).
\]
For separated rows remove the avoidance constraints touching their independent local anchor and mask inputs. Their raw means are \(R_b\le(1+e^{-n^u})\nu_{i_{z(b)}}\), whose normalized average is bounded by the fixed tag balance. The charge sum just estimated supplies the separated-product comparison. Scattered moments therefore give column sums \(O(2^n/N)=o(1)\), simultaneously, with high probability.

At such a prehistory apply Lemma~\ref{lem:clock-sampling} to the odd rows, with the predictive failures as predicates. Each predicate reads at most \(n\) rows and each row occurs in at most \(n\) predicates; the atoms and failure probabilities are smaller than any required inverse power. This gives an odd injection avoiding every predictive failure and satisfying the required joint product comparisons.

\paragraph{Even posterior rows and Hall.}
At \(v\), use the posterior \(F_x\,d\mu_{i_z}(x)/M_v\) as a row on its common neighbors. Set this row to zero off its local gate or predictive success when estimating moments. The likelihood and denominator bounds, together with \(n^{x_s}=o(na_*)\), give normalized cap
\[
 L_X\le2^n e^{-c_4na_*/2}.
\]
The same near relation has repeat cost
\(\exp(O(r\log n)-c_4na_*/2)=o(1)\),
since \(r\log n=o(na_*)\).

For separated even rows, retain the successful entering-history indicator while using the clock comparison on their disjoint odd neighbor sets. Only after replacing those outputs by their product laws do we remove nonlocal success indicators. Keep every true local star gate. Remove the anchor-avoidance constraints touching the target IDs, hold the other anchors and masks fixed, and integrate a target and its data. The cancellation is
\[
 \sum_{\mathbf y}M_v(\mathbf y)
       \frac{F_x(\mathbf y)\mu_{i_z}(x)}{M_v(\mathbf y)}
   =\mu_{i_z}(x)\sum_{\mathbf y}F_x(\mathbf y)
   \le\mu_{i_z}(x),
\]
with zero terms where \(M_v=0\). The fixed-map locality ensures that no other retained calculation consults this target ID. These integrations consequently multiply. The avoidance cost is \(1+o(1)\), and the normalized comparison means have bounded average by tag balance. Scattered moments now give even column sums \(O(2^n/N)=o(1)\), hence at most 1, with high probability on success. All even rows have mass 1 there and are supported on common neighbors of the odd injection. Hall's theorem finishes the cube and excludes both stated intermediate situations.

\end{proof}
